\section{Conclusion} \label{sec:conclusion}


% ====================================================================

%\hspace{0.6em} \textbf{Static Code Analysis} (\gls{sca}) helps developers find software defects without running the program. Traditional tools such as \textit{FindBugs} can detect common programming errors, but they can also generate false positives that developers need to check manually. As discussed by \textit{Bessey et al.} \cite{paperreference2_bessey2010few}, this becomes a bigger problem in large software projects, where \textbf{too many incorrect warnings can make developers lose trust in these tools}.

%\hspace{0.6em} The three papers analyzed in this report show different ways in which \gls{ml} can help. \textit{Kharkar et al.} \cite{paper1_kharkarlearning2022} show that Transformer based models can improve the precision of warnings generated by \textit{Infer}. \textit{Hegedűs and Ferenc} \cite{paper2_hegedHus2022static} show that simpler \gls{ml} classifiers can also give good results when trained on large datasets of real warnings. The last paper by \textit{Yerramreddy et al.} \cite{paper3_yerramreddy2023empirical} find that \gls{nns} generally achieve better results than traditional classifiers. However, their \textbf{performance depends on how the code is represented, how the data is prepared and which projects are used for testing}.

%\hspace{0.6em} Based on these research findings, \textbf{\gls{ml} can make static analysis more useful by reducing false positives and helping developers focus on warnings that are more likely to represent real defects}. There are still some challenges to address, such as the facts that an \gls{ml} model might filter out a real defect by mistake, training datasets may contain incorrect labels and models may not work as well on new projects. Another limitation is that these approaches mainly classify warnings that were already found by the static analyzer, meaning they cannot identify bugs that the original tool failed to detect.

% ====================================================================

\hspace{0.6em} The research papers analyzed in this report demonstrate that \textbf{\gls{ml} can improve traditional static analysis by making warning classification more accurate and useful for developers}. However, the effectiveness of each approach depends on the training data, source code representation and evaluation conditions.

\hspace{0.6em} The comparison also shows that \textbf{there is no universally best approach}. Simpler classifiers can provide good results with lower complexity, while neural models can capture more information from source code. Therefore, selecting an appropriate method requires considering not only classification performance, but also computational cost, generalization and the consequences of incorrectly filtering genuine defects.


% ====================================================================

\section{Future Work} \label{sec:futurework}

\hspace{0.6em} Future research could focus on \textbf{more accurate and explainable warning classification}, helping developers understand why a warning is considered a real defect or a false positive. Larger and more diverse datasets could also help \gls{ml} models work better across different projects and programming languages.

\hspace{0.6em} Another interesting direction is the use of \gls{llms} to analyze warnings together with their surrounding source code. Building on the generative approach explored by \textit{Kharkar et al.} \cite{paper1_kharkarlearning2022}, \gls{llms} could help identify possible false positives, explain why a warning might be valid and suggest code fixes. However they can also make mistakes so their results should be checked against static analysis findings or reviewed by developers. This could help reduce incorrect classifications and avoid hiding real defects. Testing such approaches in real development environments would also help understand their practical benefits.

\hspace{0.6em} From a software engineering perspective, the most promising direction is not necessarily replacing traditional static analyzers with \gls{ml} or \gls{llms}, but \textbf{combining their strengths to provide more accurate, understandable and useful results for developers}.


% ====================================================================
